% ============================================================
%  封面信息
% ============================================================
\school{国豪书院}
\major{计算机科学与技术}
\student{2452627}{徐佳涛}
\thesistitle{手写数字识别与垃圾分类}{——基于 LeNet5 的深度学习课程设计}
\thesistitleeng{Handwritten Digit Recognition and Garbage Classification}{---A LeNet5-Based Deep Learning Course Design}
\thesisadvisor{张红云\quad 教授}
\thesisdate{2026}{5}{25}

% ============================================================
%  信息说明页
% ============================================================

\infotype{thesis}

\infoabstract{%
本课程设计围绕手写数字识别与垃圾分类两个任务展开。首先在 MNIST 手写数字数据集上从零训练 LeNet5 并保存预训练权重，再将该权重迁移到 26 类垃圾分类数据集上，完成从零训练、冻结微调、全量微调、LoRA 微调和强数据增强改进实验。垃圾分类图片统一转换为灰度图，缩放到 $32\times32$，按 MNIST 均值和标准差归一化后输入 LeNet5。实验结果显示：LeNet5 在 MNIST 上达到 98.60\% 的测试准确率；垃圾分类从零训练为 3.85\%，冻结微调为 15.38\%，全量微调为 64.23\%，LoRA 微调为 25.38\%，强数据增强实验为 38.85\%。在此基础上，进一步使用 MobileNetV2 进行进阶对照，其中 MindCV 官方预训练 MobileNetV2-1.0 结合 warmup + cosine 全量微调达到 91.15\%。实验说明 MNIST 预训练权重可以作为有效初始化，但 MNIST 与垃圾分类差异较大，冻结主干适应能力有限；若追求实际准确率，保留 RGB 信息、更高输入分辨率和高质量预训练权重非常关键。
}

\infothesiswords{9000}

\infomaterials{
  \item 项目源代码（LeNet5、LoRA、数据加载、训练评估与曲线绘制模块）；
  \item 实验日志与结果文件（CSV 格式训练日志、JSON 格式结果汇总）；
  \item 训练曲线图和模板版课程设计报告 PDF。
}
